\editionPage[\pdfbookmark[0]{Corrigés des défis}{bm-defi-cor-1-parent}]{defi-cor-1}{D01 · Corrigé}{Compter jusqu’au milliard}{\hyperref[defi-ex-1]{Énoncé p.\,\pageref*{defi-ex-1}}}
\begin{blocEdition}{defi-cor-1}
\editionEtape{Une journée}$24\times3\,600=86\,400$ s.\editionEtape{Un million de secondes}$1\,000\,000=86\,400\times11+49\,600$, puis $49\,600=3\,600\times13+2\,800$ et $2\,800=60\times46+40$. Il faut $\res{11\text{ j }13\text{ h }46\text{ min }40\text{ s}}$.\editionEtape{Un milliard}Mille fois plus : environ $11\,574$ jours. Comme $11\,574=365\times31+259$, cela fait $\res{\text{plus de 31 ans}}$ sans dormir ! Les grands nombres dépassent vite l’intuition.
\end{blocEdition}
\editionPage{defi-cor-2}{D02 · Corrigé}{Entre 3,7 et 3,8}{\hyperref[defi-ex-2]{Énoncé p.\,\pageref*{defi-ex-2}}}
\begin{blocEdition}{defi-cor-2}
\editionEtape{Deux décimales}De $3{,}71$ à $3{,}79$ : $\res9$ nombres.\editionEtape{Au plus trois décimales}De $3{,}701$ à $3{,}799$ : $\res{99}$ nombres, qui comprennent les 9 précédents.\editionEtape{Tous les décimaux}$\res{\text{Impossible}}$ : avec quatre décimales on en trouve 999, avec cinq encore davantage. Entre deux décimaux différents, il y a une infinité de décimaux.
\end{blocEdition}
\editionPage{defi-cor-3}{D03 · Corrigé}{Somme 5, produit 6,09}{\hyperref[defi-ex-3]{Énoncé p.\,\pageref*{defi-ex-3}}}
\begin{blocEdition}{defi-cor-3}
\editionEtape{Premiers essais}$2{,}5\times2{,}5=6{,}25$ et $2\times3=6$. Le produit cherché, $6{,}09$, est entre les deux.\editionEtape{Affiner}$2{,}4\times2{,}6=6{,}24$ ; $2{,}3\times2{,}7=6{,}21$ ; $2{,}2\times2{,}8=6{,}16$ ; $2{,}1\times2{,}9=6{,}09$.\editionEtape{Conclure}Les nombres sont $\res{2{,}1}$ et $\res{2{,}9}$. Contrôles : $2{,}1+2{,}9=5$ et $21\times29=609$.
\end{blocEdition}
\editionPage{defi-cor-4}{D04 · Corrigé}{Les restes en escalier}{\hyperref[defi-ex-4]{Énoncé p.\,\pageref*{defi-ex-4}}}
\begin{blocEdition}{defi-cor-4}
\editionEtape{Ajouter 1}Avec 1 de plus, chaque division tombe juste : le nombre plus 1 est un multiple de 2, 3, 4, 5 et 6, donc de 60.\editionEtape{Le plus petit}$60-1=\res{59}$. Contrôle : $59=2\times29+1=3\times19+2=4\times14+3=5\times11+4=6\times9+5$.\editionEtape{Avec 7 en plus}On teste $119$, $179$, $239$… : $119=7\times17$. Le nombre cherché est $\res{119}$.
\end{blocEdition}
\editionPage{defi-cor-5}{D05 · Corrigé}{Les pains égyptiens}{\hyperref[defi-ex-5]{Énoncé p.\,\pageref*{defi-ex-5}}}
\begin{blocEdition}{defi-cor-5}
\editionEtape{Cinq pains}$\frac58=\frac48+\frac18=\res{\frac12+\frac18}$. On coupe 4 pains en deux (8 moitiés) et le dernier pain en 8 : chacun reçoit une moitié et un huitième.\editionEtape{Sept pains}$\frac78=\frac12+\frac14+\frac18$ : 4 pains en moitiés, 2 pains en quarts, 1 pain en huitièmes. Chacun reçoit $\res{\frac12+\frac14+\frac18}$.\editionEtape{Contrôle}$8\times\left(\frac12+\frac18\right)=4+1=5$ pains et $8\times\left(\frac12+\frac14+\frac18\right)=4+2+1=7$ pains.
\end{blocEdition}
\editionPage{defi-cor-6}{D06 · Corrigé}{Une somme qui se simplifie}{\hyperref[defi-ex-6]{Énoncé p.\,\pageref*{defi-ex-6}}}
\begin{blocEdition}{defi-cor-6}
\editionEtape{Trois différences}$\frac12-\frac13=\frac16$ ; $\frac13-\frac14=\frac1{12}$ ; $\frac14-\frac15=\frac1{20}$.\editionEtape{Réécrire la somme}$\left(1-\frac12\right)+\left(\frac12-\frac13\right)+\left(\frac13-\frac14\right)+\left(\frac14-\frac15\right)$.\editionEtape{Simplifier}Tout s’annule sauf le premier et le dernier terme : la somme vaut $1-\frac15=\res{\frac45}$.
\end{blocEdition}
\editionPage{defi-cor-7}{D07 · Corrigé}{Hausse puis baisse}{\hyperref[defi-ex-7]{Énoncé p.\,\pageref*{defi-ex-7}}}
\begin{blocEdition}{defi-cor-7}
\editionEtape{Hausse puis baisse}$100+10=110$ € ; 10~\% de 110 vaut 11 : $110-11=\res{99}$ €.\editionEtape{Baisse puis hausse}$100-10=90$ € ; 10~\% de 90 vaut 9 : $90+9=\res{99}$ €. Même résultat.\editionEtape{Avec 20~\%}$100\to120\to120-24=\res{96}$ €. On ne revient jamais au prix de départ : le second pourcentage porte sur un autre prix.
\end{blocEdition}
\editionPage{defi-cor-8}{D08 · Corrigé}{Les tables du banquet}{\hyperref[defi-ex-8]{Énoncé p.\,\pageref*{defi-ex-8}}}
\begin{blocEdition}{defi-cor-8}
\editionEtape{La structure}Deux places aux bouts, puis 2 places par table, une de chaque côté : nombre de places $=2+2\times$ nombre de tables.\editionEtape{Dix tables}$2+2\times10=\res{22}$ personnes.\editionEtape{Cinquante invités}$50-2=48$ et $48:2=24$ : il faut $\res{24}$ tables.
\end{blocEdition}
\editionPage{defi-cor-9}{D09 · Corrigé}{Les engrenages}{\hyperref[defi-ex-9]{Énoncé p.\,\pageref*{defi-ex-9}}}
\begin{blocEdition}{defi-cor-9}
\editionEtape{Compter les dents}15 tours de la petite roue font passer $15\times12=180$ dents.\editionEtape{Grande roue}$180:36=\res5$ tours : le nombre de tours est proportionnel au nombre de dents passées.\editionEtape{Troisième roue}Les mêmes 180 dents passent : $180:24=\res{7{,}5}$ tours.
\end{blocEdition}
\editionPage{defi-cor-10}{D10 · Corrigé}{Une ficelle autour de la Terre}{\hyperref[defi-ex-10]{Énoncé p.\,\pageref*{defi-ex-10}}}
\begin{blocEdition}{defi-cor-10}
\editionEtape{Le diamètre}Si le périmètre augmente de 1 m, le diamètre augmente de $1:\pi\approx0{,}32$ m, car le périmètre est proportionnel au diamètre.\editionEtape{La hauteur}Le diamètre s’allonge des deux côtés : la ficelle se soulève de $0{,}32:2\approx\res{16\ \mathrm{cm}}$. Un chat peut passer !\editionEtape{Le ballon}Le calcul n’utilise pas la taille de départ : $\res{\text{environ 16 cm aussi}}$. Résultat surprenant, mais juste.
\end{blocEdition}
\editionPage{defi-cor-11}{D11 · Corrigé}{Le carré penché}{\hyperref[defi-ex-11]{Énoncé p.\,\pageref*{defi-ex-11}}}
\begin{blocEdition}{defi-cor-11}
\editionEtape{Le grand carré}Aire : $4\times4=16$ cm$^2$.\editionEtape{Les quatre coins}Les quatre triangles, assemblés deux par deux, forment deux rectangles de 1 cm sur 3 cm : $2\times3=6$ cm$^2$.\editionEtape{Le carré penché}$16-6=\res{10\ \mathrm{cm}^2}$. On peut contrôler en comptant les carreaux entiers et les morceaux.\label{fig-defi-cor-11}\figureGeo[corrige]{defi-carre}

\end{blocEdition}
\editionPage{defi-cor-12}{D12 · Corrigé}{Les deux feux}{\hyperref[defi-ex-12]{Énoncé p.\,\pageref*{defi-ex-12}}}
\begin{blocEdition}{defi-cor-12}
\editionEtape{En secondes}45 s et 1 min 20 s $=80$ s.\editionEtape{Multiples}Multiples de 80 : 80, 160, 240, … 720. Or $720=45\times16$ : c’est le premier multiple commun.\editionEtape{Heure}$720$ s $=12$ min : les feux clignotent ensemble à $\res{8\text{ h }12\text{ min}}$.
\end{blocEdition}
\editionPage{defi-cor-13}{D13 · Corrigé}{Combien de points ?}{\hyperref[defi-ex-13]{Énoncé p.\,\pageref*{defi-ex-13}}}
\begin{blocEdition}{defi-cor-13}
\editionEtape{Petit rayon}$r=1$ : le cercle de centre B est trop petit, $\res0$ point. $r=2$ : les cercles se touchent, $\res1$ point.\editionEtape{Rayon moyen}$r=5$ : les cercles se coupent en $\res2$ points.\editionEtape{Grand rayon}$r=10=4+6$ : ils se touchent de nouveau, $\res1$ point. $r=11$ : le cercle de centre A est entièrement à l’intérieur de l’autre, $\res0$ point.\label{fig-defi-cor-13}\figureGeo[corrige]{defi-cercles}

\end{blocEdition}
\editionPage{defi-cor-14}{D14 · Corrigé}{Le centre perdu}{\hyperref[defi-ex-14]{Énoncé p.\,\pageref*{defi-ex-14}}}
\begin{blocEdition}{defi-cor-14}
\editionEtape{Construire}On trace deux cordes [AB] et [CD], puis leurs médiatrices au compas. Elles se coupent en O, en vert.\editionEtape{Justifier}Le centre est à égale distance de A et de B (deux rayons) : il est sur la médiatrice de [AB]. De même, il est sur la médiatrice de [CD].\editionEtape{Conclure}Le centre appartient aux deux médiatrices : c’est leur point d’intersection $\res O$.\label{fig-defi-cor-14}\figureGeo[corrige]{defi-centre}

\end{blocEdition}
\editionPage{defi-cor-15}{D15 · Corrigé}{Les aiguilles en mouvement}{\hyperref[defi-ex-15]{Énoncé p.\,\pageref*{defi-ex-15}}}
\begin{blocEdition}{defi-cor-15}
\editionEtape{Vitesses}La grande aiguille tourne de $6^\circ$ par minute, la petite de $0{,}5^\circ$ par minute.\editionEtape{2 h 30}Petite : $2\times30+30\times0{,}5=75^\circ$. Grande : $30\times6=180^\circ$. Angle : $180-75=\res{105^\circ}$.\editionEtape{9 h 45}Petite : $9\times30+45\times0{,}5=292{,}5^\circ$. Grande : $45\times6=270^\circ$. Angle : $\res{22{,}5^\circ}$.
\end{blocEdition}
\editionPage{defi-cor-16}{D16 · Corrigé}{Des triangles d’allumettes}{\hyperref[defi-ex-16]{Énoncé p.\,\pageref*{defi-ex-16}}}
\begin{blocEdition}{defi-cor-16}
\editionEtape{Lister}Côtés possibles : (2, 5, 5) ; (3, 4, 5) ; (4, 4, 4). Les autres essais, comme (3, 3, 6), ne ferment pas : $3+3=6$, le triangle est aplati.\editionEtape{Compter}$\res3$ triangles différents.\editionEtape{Classer}(2, 5, 5) est isocèle ; (4, 4, 4) est équilatéral, donc aussi isocèle ; (3, 4, 5) n’a aucun côté égal : il est même rectangle.
\end{blocEdition}
\editionPage{defi-cor-17}{D17 · Corrigé}{Deux miroirs parallèles}{\hyperref[defi-ex-17]{Énoncé p.\,\pageref*{defi-ex-17}}}
\begin{blocEdition}{defi-cor-17}
\editionEtape{Un exemple}M est à 1 carreau de $(d_1)$ : $M_1$ est 1 carreau de l’autre côté, puis $M_2$ est obtenu en traversant $(d_2)$. Sur la figure, $MM_2=\res{6\ \mathrm{cm}}$.\editionEtape{Un autre point}On retrouve $MM_2=6$ cm, toujours perpendiculairement aux droites.\editionEtape{Bilan}Deux symétries successives d’axes parallèles reviennent à faire glisser la figure de $\res{2\times3=6\ \mathrm{cm}}$ : on retrouve l’orientation de départ.\label{fig-defi-cor-17}\figureGeo[corrige]{defi-miroirs}

\end{blocEdition}
\editionPage{defi-cor-18}{D18 · Corrigé}{La pyramide de cubes}{\hyperref[defi-ex-18]{Énoncé p.\,\pageref*{defi-ex-18}}}
\begin{blocEdition}{defi-cor-18}
\editionEtape{Quatre étages}$1+4+9+16=\res{30\ \mathrm{cm}^3}$.\editionEtape{Cinq étages}$30+25=\res{55\ \mathrm{cm}^3}$.\editionEtape{Dépasser 200}$55+36=91$ ; $91+49=140$ ; $140+64=204$ : il faut $\res8$ étages.\label{fig-defi-cor-18}\figureGeo[corrige]{defi-pyramide}

\end{blocEdition}
\editionPage{defi-cor-19}{D19 · Corrigé}{Le tableau à reconstituer}{\hyperref[defi-ex-19]{Énoncé p.\,\pageref*{defi-ex-19}}}
\begin{blocEdition}{defi-cor-19}
\editionEtape{Garçons}$50-28=22$ garçons, dont 12 à la cantine et 10 qui n’y déjeunent pas.\editionEtape{Filles}À la cantine : $30-12=\res{18}$ filles ; hors cantine : $28-18=\res{10}$ filles.\editionEtape{Proportion}$\frac{18}{28}=\frac9{14}$ des filles, soit environ $\res{64\text{~\%}}$ (valeur arrondie).
\end{blocEdition}
\editionPage{defi-cor-20}{D20 · Corrigé}{Les chaussettes dans le noir}{\hyperref[defi-ex-20]{Énoncé p.\,\pageref*{defi-ex-20}}}
\begin{blocEdition}{defi-cor-20}
\editionEtape{Une paire}Avec 2 chaussettes, la malchance donne une rouge et une bleue. La troisième est forcément de l’une des deux couleurs : $\res3$ chaussettes suffisent, et la probabilité d’avoir une paire vaut alors 1.\editionEtape{Deux rouges}Au pire, on tire d’abord les 5 bleues, puis une rouge : il faut $\res7$ chaussettes pour être sûr d’en avoir deux rouges.\editionEtape{Retenir}« Être certain » correspond à une probabilité de 1 : on raisonne sur le cas le plus défavorable.
\end{blocEdition}
\editionPage{defi-cor-21}{D21 · Corrigé}{La spirale du robot}{\hyperref[defi-ex-21]{Énoncé p.\,\pageref*{defi-ex-21}}}
\begin{blocEdition}{defi-cor-21}
\editionEtape{Distance}$1+1+2+2+3+3+4+4+5+5=\res{30}$ cases.\editionEtape{Position}Horizontalement : $+1-2+3-4+5=3$. Verticalement : $+1-2+3-4+5=3$. Le robot est à $\res{\text{3 cases à droite et 3 cases au-dessus de O}}$.\editionEtape{Avec des répétitions}On peut écrire : « répéter 2 fois : avancer de 1, tourner à gauche », puis la même chose avec 2, 3, 4 et 5 : la spirale apparaît en vert.\label{fig-defi-cor-21}\figureGeo[corrige]{defi-spirale}

\end{blocEdition}
